\section{The joint acquired-geometry transfer theorem}\label{sec:transfer}
The following assumptions are imposed before choosing the observer or a codebook. They use the actual one-step laws in \eqref{eq:Lambda}.

\begin{definition}[Clock neutrality and block-stable acquisition]\label{def:hypotheses}
There are constants $\rho<\infty$, $B<\infty$, an integer $b\ge1$, $0<\theta<1$, and positive numbers $\beta_t\le B$, such that:

\textup{(i)} For all $t,s,a$, the report law $J_t(s,a)$ is equivalent to a fixed probability reference measure $m^a$. The same reference is used at every phase. Thus a report does not reveal a previously uncharged clock through disjoint supports.

\textup{(ii)} For all $t,s,a$, $\Lambda_t(s,a)$ has density at most $\rho$ relative to $d$-dimensional Lebesgue measure on $K$.

\textup{(iii)} For all $t,s,s',a$,
\begin{equation}\label{eq:stability}
 W_1(\Lambda_t(s,a),\Lambda_t(s',a))\le\beta_t\norm{s-s'},
 \qquad \prod_{j=t}^{t+b-1}\beta_j\le\theta
\end{equation}
whenever the indicated block is present. The report kernels are jointly Borel and continuous in total variation in $s$, uniformly over the finite action family.
\end{definition}
For a sequence of horizons, the same constants are required. The last incomplete block is controlled by $B$. The definition permits $\beta_t>1$; it does not assume pathwise contraction of two filters on every report. The raw-kernel criteria in Section~\ref{sec:kernels} verify the report continuity and Wasserstein property directly. A density bound is stronger than positive noise, but no compact interior carrier or lower bound on acquired density is required.

\begin{theorem}[Charged-clock acquired-geometry transfer]\label{thm:main}
Consider a prepared experiment with Definition~\ref{def:affine}, the task and machine model of Section~\ref{sec:problem}, and Definition~\ref{def:hypotheses}. There exist $0<c\le C<\infty$, depending only on $K$, its norm, $d,\rho,B,b,\theta,L_g,\alpha$, such that for every $n\ge1$ and every $M\ge1$,
\begin{equation}\label{eq:maincp}
 cM^{-2/d}\le R_\cp(n,M)-B_n^*\le CM^{-2/d}.
\end{equation}
The exact-$n$ autonomous class is nonempty precisely when $M\ge n+1$. For every such $M$,
\begin{equation}\label{eq:mainaut}
 c(M-n)^{-2/d}\le R_\aut(n,M)-B_n^*\le C(M-n)^{-2/d}.
\end{equation}
The constants may be decreased or enlarged once so that the same pair works in both bounds. The upper is realized by a deterministic observer with at most $M$ total labels, including all controller phases. It uses no discarded history or uncharged exact state.

More generally, without the block-product restriction, put
\begin{equation}\label{eq:weights}
 a_t=L_g\prod_{j=t}^{n-1}\beta_j,\quad
 \gamma=\frac{d}{d+2},\quad A_n=\sum_{t=1}^n a_t^\gamma.
\end{equation}
The empty product is one. There is a constant $C_0$ depending only on $K,d,\rho$ and the norm comparison such that
\begin{equation}\label{eq:weighted}
 R_\aut(n,M)-B_n^*\le C_0 A_n^{1/\gamma}(M-n)^{-2/d},\qquad M\ge n+1.
\end{equation}
Thus the continuation-sensitivity profile, not equal phase allocation, determines the upper. Block stability makes $A_n$ uniformly bounded.
\end{theorem}

The rate is attached to the dimension of an actually acquired law after a real report. The theorem does not replace an unavailable acquired law by the formal dimension of $K$: condition (ii) forces a full-dimensional probability distribution at each call. Section~\ref{sec:strata} replaces $d$ by the largest actually acquired stratum dimension under explicit mass assumptions.

\begin{corollary}[Resolution and nonquadratic tasks]\label{cor:resolution}
For target excess $\eta>0$ below a fixed constant, the required checkpoint labels are of order $\eta^{-d/2}$ and the required autonomous labels are
\[
 n+\Theta(\eta^{-d/2}).
\]
Persistent bits are therefore $\lceil\log_2(n+\Theta(\eta^{-d/2}))\rceil$, not programme length or total bit complexity.

On a finite terminal audit alphabet, let $\Phi$ be a $C^1$ strongly convex function on its probability simplex and let
\[
 \ell(j,u)=\Phi(e_j)-\Phi(u)-\nabla\Phi(u)\cdot(e_j-u).
\]
If the audit probability map $P:K\to\Delta$ is affine and injective, these bounded proper losses satisfy the task hypotheses, after harmless rescaling. For example $\Phi(u)=\sum_j u_j^2+\zeta\sum_j u_j^4$, $\zeta>0$, is nonquadratic. The theorem applies with constants depending on $P$ and $\Phi$.
\end{corollary}
\begin{proof}
The first assertion is the inversion of \eqref{eq:maincp}--\eqref{eq:mainaut}. Conditional expectation of the displayed loss is uniquely minimized by the audit distribution, and its Bayes risk is $\sum_jP_j(s)\Phi(e_j)-\Phi(P(s))$. Strong convexity and injectivity give strong concavity in $s$; compactness and $C^1$ regularity give bounded loss and a Lipschitz Bayes risk. This is an example class, not a claim for every bounded loss.
\end{proof}

\subsection{An actual-mass and continuation-curvature bound}
For a concave value $v$ on $K$, a finite acquired measure $\lambda$ and a partition $\mathcal P$, denote the nonnegative expression in \eqref{eq:taskdist} by $\mathcal J_v(\lambda,\mathcal P)$. It is invariant under affine changes of predictive coordinates. Lemma~\ref{lem:curvature} supplies a convex Lipschitz extension $f$ of $-v$ to all of $\R^d$. Let $\kappa_f=\Delta f$ be its positive distributional curvature measure. For a grid of side $h$, with cell centers $z_C$, define
\begin{equation}\label{eq:masscurvature}
 \mathfrak C_h(\lambda,f)=h^{2-d}\sum_C\lambda(C)
 \kappa_f\bigl(\overline B_{3\sqrt d h}(z_C)\bigr).
\end{equation}
Only cells meeting the compact carrier occur. The estimate
\[
 \mathcal J_v(\lambda,\mathcal P_h)\le C_d\mathfrak C_h(\lambda,f)
\]
holds also for singular measures. It does not normalize their mass or smooth a switching set. A bound on acquired cell mass and a bound on continuation curvature are separate inputs. Under (ii), bounded overlap reduces this expression to $C\rho\Lip(v)h^2$. Equation~\eqref{eq:masscurvature} is a useful sufficient local quantity, not a necessary-and-sufficient characterization or an assertion that every singular law has exponent $2/d$.
